\section{Limitations and Implications}

\paragraph{What is established}
The stopping theorem covers three specified ideal mechanisms on finite-state
Markov sources with positive varentropy and bounded token probabilities. The
verified-rule proof uses exact reals and continuous rotations; its finite
numerical implementation is tested separately. The current-upstream study is
GPT-2 only. The earlier GPT-2/Qwen studies use a different, explicitly processed
source and adapter. These complementary results do not constitute a theorem
for arbitrary dependent language-model distributions.

\paragraph{Observer and sample limits}
The observer knows the prompt, source, and encoded length. Unknown prompts,
unknown-length mixtures, model mismatch, and natural tails require separate
evaluation. The upstream design has ten held-out prompt clusters, with nine
represented in the eligible 512-bit stopping slice; the older controlled
studies have five held-out topic families. Canonical-text selection is
substantial in some cohorts. All outcomes remain available, but conditional
AUCs are not unconditional transmitted-text security statements. The 128-bit
upstream FPR of $0.200$ also shows why high AUC does not imply a low-FPR
operating guarantee. Endpoint replay is a small post-hoc no-tail diagnostic,
with known plaintext length and a structural true-positive property.

\paragraph{Coding and delivery assumptions}
The guarded theorem assumes uniform payloads with independent guard bits and
uses the same integer-state stopping reference. Authentication does not prove
that the entire framed packet is pseudorandom. Public endpoint support can
persist for arbitrary delivered packet bytes, even when that sampling premise
is unavailable. Numerical replay requires more than matching weights: the
retained thread-count mismatch changes integer frequencies. Cross-hardware
identity is not claimed. Weak passphrases can be guessed, and active edits
can prevent recovery; authentication is not visible-text integrity.

\paragraph{Design implications}
Guard bits control a concrete precision defect while leaving the stopping
policy observable. Reverse verification repairs a decoding failure but does
not remove terminal-information concentration in the ideal rule or the
measured signal in this upstream cohort. Exogenous endpoints decouple length
from their own content, with a delivery cost that can be traded for extra
independently budgeted tokens. Retrying until recovery changes that comparison.
A content-dependent stopping rule shared with the declared natural source
is another reference choice. The source law, failure population, and endpoint
must therefore be specified together; token-table fidelity alone does not
settle complete-text security.
